\chapter{The Defeat of the Opposition}

The period between the fourteenth party congress in December 1925 and the fifteenth in December 1927, which was marked by the beginnings of effective planning, the first steps in a programme of intensive industrialization, and the ripening of the agricultural crisis, also saw the culmination of an embittered party conflict, in which Trotsky fiercely and unsuccessfully challenged Stalin's growing monopoly of power. When the triumvirate broke up, and Stalin defeated his rivals at the fourteenth congress, Trotsky remained haughtily silent; Zinoviev and Kamenev had in the past equalled, and sometimes outdone, Stalin in the vehemence of their attacks on him. But, when Zinoviev and Kamenev took up the cause of industrialization against the peasant orientation of Stalin and Bukharin, and when Stalin's personal ambitions became more open and more menacing, neutrality was no longer possible. In the summer of 1926, Trotsky, Zinoviev and Kamenev with their followers constituted themselves as a ``united opposition'', and appeared as such at the July session of the party central committee. The sequel showed the strength of Stalin's control of the party machine. Trotsky was at first cautiously handled. But Zinoviev lost his seat in the Politburo, and Kamenev his government offices. The united opposition enjoyed much sympathy in the party ranks. But its active supporters were not more than a few thousand; and these on various pretexts were constantly harassed by the authorities. 
% 注：去除了 "^possible." 中的标记，并修复了 "unsuc- cessfully" 的断字。

Lack of internal cohesion and of mutual confidence added to the handicaps of the opposition, which took the field without any clear-cut line other than denunciation of the party leaders. The public withdrawal of the mutual accusations which Trotsky, Zinoviev and Kamenev had levelled at one another in the past three years invited ridicule. Zinoviev's vacillating temperament and proneness to compromise made him uncongenial to Trotsky, who, having cast off earlier inhibitions, called for an unflinching offensive against Stalin. Scarcely had the united opposition been formed when, by an unfortunate coincidence, the New York Times published for the first time the text of Lenin's testament. Though Trotsky was certainly not privy to the publication, the supposition that knowledge of the document had come originally from him, or from sources close to him, was not unfounded. Bitterness between the two protagonists in the struggle reached its peak. At a heated session of the Politburo Trotsky branded Stalin as ``the grave-digger of the revolution''; and the party central committee, reacting to the growing tension, dismissed Trotsky from his membership of the Politburo. At a party conference in October 1926 and a session of IKKI a month later, Stalin attacked Trotsky in language of increasing virulence, vindictively exhuming his pre-1914 record of flirtation with Menshevism and acrimonious exchanges with Lenin. The opposition was indicted not only for ``fractionalism''---a sin condemned by the party congress of 1921---but for a ``social-democratic deviation''. Stalin was, however, still content to bide his time, and did not bring the issue to breaking-point. 
% 注：去除了 "I Lack" 的边缘噪点，去除了 ".At a heated" 的前导句号，修复了 "social-de- mocratic" 的跨行断字。

In the spring of 1927 the turn of events in China spurred Trotsky to fresh protests; and in May the opposition issued a document, drafted mainly by Trotsky, and known from the number of its original signatories as ``the declaration of the 83'', which offered the fullest expose hitherto available of its views. Apart from its excursion into foreign affairs, the declaration denounced current agricultural policy for ignoring the process of ``differentiation'' among the peasantry, and for neglecting the poor peasant in order to bolster up the kulak. In general terms, it accused the party leaders of substituting ``the petty-bourgeois theory of socialism in one country'' for ``Marxist analysis of the real situation of the proletarian dictatorship in the USSR'', and of favouring ``Rightist, non-proletarian and anti-proletarian elements'' inside and outside the party. It also demanded full publicity for the views of the opposition. Coming at a moment when the leaders had been severely shaken by Chiang Kai-shek's volte-face in China and by the rupture of relations with Britain, this was a telling blow. A month later, on flimsy pretexts, Trotsky and Zinoviev were summoned before the party control commission, the organ entrusted with the maintenance of party discipline, and threatened with expulsion from the party. After angry exchanges, the question was referred to the party central committee where the battle continued, both Trotsky and Stalin speaking more than once. The only novel feature of the debate was the charge against Trotsky of disloyalty to the Soviet state in face of its enemies. He was now branded not merely as a heretic, but as a traitor (``the united front from Chamberlain to Trotsky''). Finally the opposition was induced to sign a declaration re-affirming its unconditional loyalty to the national defence of the USSR, and disclaiming any desire to split the party or found a new party. On these terms the proposal to expel Trotsky and Zinoviev was shelved. 
% 注：修正了 "Alter angry exchanges" (由于 OCR 错误将 f 识别成了 l，改为 After)。

This respite did not, however, mean any remission in the hounding of the opposition. The struggle against Trotsky was the occasion for the introduction or perfection of many of the instruments of control characteristic of Stalin's dictatorship. Since the first attacks on Trotsky at the end of 1924, the access of the opposition to the press had been severely restricted. Zinoviev had been muzzled when his Leningradskaya Pravda was taken over in January 1926. Now the ban became absolute. Articles on the Chinese crisis submitted by Trotsky in April 1927 to Pravda and Bol'shevik were rejected. Throughout the summer attacks of increasing violence on him and his supporters were published in the press without any right of reply. Opposition meetings were interrupted and broken up by hooligans. The opposition submitted to the party central committee an extensive ``platform'' of its views, once more drafted mainly by Trotsky, and demanded that it should be printed and circulated in preparation for the party congress, now fixed for December 1927; this was refused. Attempts were made to print it illegally. On September 12 the OGPU discovered the illicit press, and arrested those engaged in the work. Fourteen members of the party were expelled; and Preobrazhensky, who admitted complicity, was added to the number. The occasion was significantly remembered as the first time when the police power of the OGPU had been invoked to quell dissent in the party. 
% 注：去除了 "[the occasion" 和 "(summer" 的单侧括号。

From this point events moved steadily to their predestined conclusion. Mass meetings were organized to denounce the opposition and demand the expulsion of its leaders. Well-known supporters of the opposition were removed from the scene of action by appointments to posts in distant parts of the USSR or to diplomatic posts abroad. At a session of the presidium of IKKI on September 29, Trotsky delivered a two-hour indictment of Stalin's policies. He was then voted out of IKKI with only two dissentients. The same experience was repeated a month later, amid scenes of violence, at the party central committee. Stalin himself proposed the removal of Trotsky and Zinoviev from the committee; this was carried, apparently without a vote. The Moscow police were active at the tenth anniversary celebrations of the revolution on November 7, harassing Trotsky and other opposition leaders as they drove round the city, and seizing banners with opposition slogans. Zinoviev met similar treatment in Leningrad. These public appearances of the opposition leaders were denounced as hostile demonstrations. A week later Trotsky and Zinoviev were expelled from the party by vote of the party central committee; and Kamenev and several others were dropped from the committee. 

When therefore the party congress assembled in December 1927, Trotsky and Zinoviev were absent, and the occasion was something of an anti-climax. Twelve opposition members were removed from the party central committee. Kamenev and Rakovsky, who made the principal speeches for the opposition, were frequently interrupted; and their case was weakened by would-be conciliatory approaches behind the scenes to the party leaders, which were contemptuously repulsed. The congress expelled from the party 75 ``active workers of the Trotskyite opposition'' and 15 other dissidents. Trotsky and Zinoviev were replaced in the Politburo by Kuibyshev and Rudzutak, both stalwarts of the official line. But Trotsky, though expelled, had not been silenced and was still dangerous. The Politburo decided to remove him and about 30 of his chief supporters from Moscow. Most of these were assigned to minor official posts in Siberia or Central Asia. Trotsky refused such an assignment, and was forcibly deported under an article of the criminal code relating to counter-revolutionary activities. Exceptionally---for they were recognized as presenting no danger---Zinoviev and Kamenev were banished to Kaluga, only a few hundred miles from Moscow; and even this sentence was not rigidly enforced. Trotsky's place of exile was Alma-Ata, a town on the furthest confines of Soviet Central Asia, remote even from the railway. Here he remained till his deportation from the USSR a year later, conducting a slow, but voluminous, correspondence with members of the opposition scattered all over Siberia, receiving from time to time secret reports from supporters still at large in Moscow, and addressing a flow of protests, personal and political, to the authorities. The defeat of the united opposition, and the expulsion of the only figure in the party whose stature made him a match for Stalin, was a historical landmark. 
% 注：去除了 "^Trotskyite" 和 "(refused" 等乱码。极严重跨行错乱修复：将 "iThe defeat of the united opposition, and the expulsion to the authorities. of the only figure..." 拆解拼合，将 "to the authorities" 放回上一句的结尾，将 "The defeat..." 还原为独立句首。

When the party congress of 1921 banned ``fractionalism'' and the propagation of dissentient opinion, the aim was to maintain the unity of the party and the loyalty of its members. Dissent in the party carried party sanctions, but was not yet disloyalty to the state. Party representatives of state institutions were obliged to follow the party line and to speak with a single voice. But this obligation did not extend to non-party representatives. By 1927 the distinction between party and state had been gradually eroded. Economic as well as political emergencies strengthened the need for a firm undivided authority. ``The supreme historical task of building a socialist society'', declared a resolution of the party conference in October 1926, ``imperatively demands a concentration of the forces of party, state and working class on questions of economic policy''. Decrees were now sometimes issued in the joint name of the party central committee and the central executive committee of the Congress of Soviets. The power of the state was available both to enforce party edicts and to enforce party discipline on party members. Supreme authority in party and state was concentrated in one institution---the Politburo of the party; and that authority was absolute. It was significant that the opposition headed by Trotsky was the last to be officially designated by that name; the word, familiar in the practice of western democracy, implied opposition to the ruling party which was not incompatible with loyalty to the state. In the next stage, dissent was described as ``deviation''---the language not of political differences, but of doctrinal heresy. Finally, dissident groups were simply branded as ``anti-party'', hostility to the party being unconditionally identified with hostility to the state. 
% 注：去除了 "P926" (修正为 1926), "'of party", "liparty discipline", "—^the /Politburo" 等严重识别噪点。

The elimination of legal opposition was part of a process which concentrated and centralized the combined authority of party and state, and made it absolute. The results were often unpremeditated, but were none the less irresistible. The same forces were at work in many fields. The limited licence hitherto accorded in the press and in journals to expressions of independent opinion on peripheral issues (sometimes accompanied by an editorial reservation) now almost entirely disappeared, control being silently achieved not by direct censorship, but by changes in editors and editorial boards. A proliferation of different literary schools---some avant-garde, some formalist, some professedly proletarian---had been characteristic of the years after the revolution. A pronouncement of the party central committee in 1925, apparently drafted or inspired by Bukharin, showed a willingness to tolerate these multifarious approaches to literature, none of them directed against the regime, and a reluctance to choose between them. Among these literary organizations was one calling itself the All-Russian Association of Proletarian Writers (VAPP), dominated by an ambitious literary politician named Averbakh, who had good party connexions, and from 1926 onwards conducted a campaign, in the name of a ``cultural revolution'', to place VAPP in control of all literary productions, and to suppress the publications of other schools. It was not till December 1928, after prolonged resistance, that the party central committee issued a decree placing all publishing under party and state control, which was in practice exercised through VAPP. It seems clear that this consummation had not been planned, or even perhaps desired, by the central committee, and least of all by Stalin. But corruption spread from the top. Petty dictators at lower levels eliminated their rivals by cajoling and flattering the superior authority, and by imitating its methods. 

The drive to strengthen and centralize authority was especially conspicuous in the domain of law. The administration of law had originally been reserved to the constituent republics of the USSR, each of which had its own courts and its own People's Commissariat of Justice. But the constitution of the USSR of 1923 provided for a Supreme Court of the USSR with powers to rule on questions of law submitted to it by the Supreme Courts of the republics; and the Praesidium of the TsIK appointed a Procurator whose function it was to supervise the administration of law throughout the USSR. The constitution also established a Unified State Political Administration (OGPU)---the heir of the original Cheka, whose name it often continued to bear in popular speech---to control the GPUs of the republics, which now became local agencies of a powerful central organ. While each republic had its own criminal code (the code of the RSFSR serving in practice as a model for the rest), the USSR issued in 1924 a set of ``Foundations of Criminal Legislation'', which sought to reserve for the exclusive competence of the USSR ``state crimes'', alternatively described as ``counter-revolutionary crimes'', and crimes threatening ``the administrative order''. The republics were instructed to bring their criminal codes into line with the ``foundations''. The task was evidently approached with reluctance. It was not completed for the RSFSR till the middle of 1927, and somewhat later for the other republics. 
% 注：去除了 "(of the USSR" 的括号。

Centralization of authority was accompanied by a gradual modification of current attitudes to law. The Marxist conception of law as an instrument of class rule, destined eventually to wither away with the state, and in the meanwhile to be administered with special indulgence for workers and peasants, was silently abandoned. The market practices of NEP demanded the development and strict enforcement of civil law. The maintenance of law and order under the label of ``revolutionary legality'' became a major objective. Initial emphasis on the reformatory rather than the punitive aspects of penal policy faded away. These changes reflected the growing economic and political tension. Such episodes as the assassination of the Soviet representative in Warsaw in June 1927, and the explosion of a bomb in Leningrad a few days later, led to a loud outcry against monarchists, provocateurs and agents of foreign governments; and the demand for what were officially called ``measures of social defence'' automatically increased the powers and prestige of the OGPU. The tenth anniversary of the creation of the Cheka was enthusiastically celebrated in December 1927, a few weeks after the tenth anniversary of the revolution. In March 1928 an instruction ``On Penal Policy and the Regime of Places of Confinement'' paved the way for an extension of the hitherto limited network of ``concentration camps'' for political offenders under the management of the OGPU, and prescribed the severest measures of repression for ``dissidents and professional criminals and recidivists''. The year 1928, which followed the defeat of the opposition, and was marked by the growing pressures of industrialization, witnessed throughout Soviet society the imposition from above of a powerful and despotic authority, of a rigid orthodoxy of opinion, and of the harshest penalties on those who offended against it. 
% 注：去除了段首的逗号，修复了断词 "pea- I'sants" 为 "peasants"。极严重跨行错位修复：解开 "provocateurs and (were officially called “measures of social defence” automati- agents of foreign governments; and the demand for what cally increased" 的混乱，重新梳理为通顺的两句结构。